\section*{Hoofdstuk \ref{ch:b3:virology} --- Virologie}

\begin{solution}{exo:b3:virology:1}
Polio: IV, RNA van de positieve streng --- een RNA-afhankelijk
RNA-polymerase. Griep: V, RNA van de negatieve streng --- een eigen
RNA-polymerase, meegedragen in het \omterm{def:b3:virology:virus}{virion}. Hiv: VI --- reverse
transcriptase (en integrase). Herpes simplex: I, dubbelstrengs DNA ---
de gastheer zou in beginsel alles kunnen leveren, maar het \omterm{def:b3:virology:virus}{virus} brengt
een eigen DNA-polymerase en thymidinekinase mee. Hepatitis B: VII ---
een reverse transcriptase om zijn pregenomisch RNA in DNA te
kopiëren. Rotavirus: III, dubbelstrengs RNA --- een RNA-afhankelijk
RNA-polymerase in het deeltje.
\end{solution}

\begin{solution}{exo:b3:virology:2}
Het DNA van de faag ($^{32}$P) ging de bacteriën binnen en het eiwit
($^{35}$S) bleef erbuiten en kon eraf worden geslagen; het nageslacht
van de faag erfde het $^{32}$P. Het DNA is dus wat de faag inspuit en
wat het maken van nieuwe fagen stuurt. Had het eiwit de instructies
gedragen, dan zou het $^{35}$S binnen zijn gevonden en aan het
nageslacht zijn doorgegeven.
\end{solution}

\begin{solution}{exo:b3:virology:3}
$84/(0.1\times 10^{-7}) = 8.4\times 10^{9}$ plaquevormende eenheden per
mL. De microscoop telt elk deeltje, ook lege \omterm{def:b3:virology:virus}{capsiden}, deeltjes met
gebrekkige genomen en deeltjes die er niet in slagen te hechten of in te
spuiten; alleen deeltjes die een besmetting voltooien maken plaques.
\end{solution}

\begin{solution}{exo:b3:virology:4}
Hechting: intrederemmers (maraviroc), neutraliserende antilichamen.
Intrede en fusie: fusieremmers, blokkers van de aanzuring van het
\omterm{def:b3:membrane-traffic:endocytosis}{endosoom}. Ontmanteling: middelen die het \omterm{def:b3:virology:virus}{capside} binden (pleconaril;
lenacapavir). Expressie en replicatie: nucleoside-analogen, remmers van
het reverse transcriptase en van het integrase, interferon,
\omterm{def:b3:rna-regulation:rnai}{RNA-interferentie}. Assemblage: \omterm{def:b3:virology:drugs}{proteaseremmers} (ongeknipte
polyproteïnen kunnen niet worden samengesteld). Vrijlating:
neuraminidaseremmers (griep), en cytotoxische T-cellen die de cel doden
voordat zij vrijlaat.
\end{solution}

\begin{solution}{exo:b3:virology:5}
Eén eiwit van \qty{30}{kDa} is $30\,000/110 = 273$ residuen, $818$
nucleotiden: \qty{0.82}{kb}, \qty{18}{\%} van een genoom van
\qty{4.5}{kb}. De $180$ kopieën elk afzonderlijk coderen zou
\qty{147}{kb} kosten, drieëndertig keer het hele genoom; het
capside-eiwit zelf (\qty{5.4}{MDa}) weegt viermaal zoveel als het
genoom (\qty{1.5}{MDa}).
\end{solution}

\begin{solution}{exo:b3:virology:6}
$V = 10^{5}(3e^{-0.5t} - 0.5e^{-3t})/2.5$: $t = 0.5$: $8.9\times
10^{4}$; $t = 2$: $4.4\times 10^{4}$; $t = 10$: $8\times 10^{2}$. Onder
$50$ wanneer $1.2\times 10^{5}e^{-0.5t} = 50$: $t = 2\ln(2400) \approx
\qty{16}{d}$.
\end{solution}

\begin{solution}{exo:b3:virology:7}
$uL = 0.36$; foutloze fractie $e^{-0.36} = 0.70$. Verdrievoudigd: $uL =
1.08$, $e^{-1.08} = 0.34$ --- het meeste nageslacht draagt nu
mutaties en de fitste sequentie wordt naar de drempel toe verdund. Een
coronavirus van \qty{30}{kb} met $u = 10^{-4}$ zou $uL = 3$ hebben en
maar \qty{5}{\%} foutloze kopieën, voorbij de drempel; zijn exonuclease
brengt $u$ op $10^{-5}$, $uL = 0.3$, drie kwart foutloos.
\end{solution}

\begin{solution}{exo:b3:virology:8}
Drift: puntmutaties in het hemagglutinine en het neuraminidase hopen
zich onder de selectie door antilichamen op, zodat de stam van elk
seizoen deels nieuw is en de afweer van vorig jaar deels verouderd.
Shift: een gelijktijdige besmetting van één cel door een menselijke en
een dierlijke stam herschikt de acht segmenten, en er verschijnt een
hemagglutininesubtype van vogels of varkens in een \omterm{def:b3:virology:virus}{virus} dat aan de
mens is aangepast. Niemand heeft antilichamen tegen het nieuwe subtype,
zodat de hele bevolking in één keer vatbaar is --- de voorwaarde voor
een pandemie, waaraan in 1918, 1957, 1968 en 2009 telkens door een
nieuw hemagglutinine werd voldaan.
\end{solution}

\begin{solution}{exo:b3:virology:9}
Het thymidinekinase van het herpesvirus fosforyleert aciclovir, dat de
cellulaire kinasen nauwelijks aanraken; de cellulaire kinasen maken
daarna het trifosfaat af, dat het virale DNA-polymerase inbouwt en
dat, zonder 3$'$-hydroxylgroep, de keten beëindigt. Onbesmette cellen
maken de werkzame vorm nooit, en het virale polymerase verkiest haar
boven het cellulaire. Resistentie: verlies of verandering van het virale
thymidinekinase (het meest voorkomend), of mutaties in het virale
polymerase die het analoog uitsluiten.
\end{solution}

\begin{solution}{exo:b3:virology:10}
Met $T$ constant zijn de vergelijkingen lineair in $(I, V)$ met matrix
$\begin{pmatrix} -\delta & \beta T \\ p & -c \end{pmatrix}$, spoor
$-(\delta + c) < 0$ en determinant $\delta c - \beta Tp$. De
eigenwaarden hebben een negatief spoor, dus is er precies één positief
wanneer de determinant negatief is: $\beta Tp > c\delta$, dat wil
zeggen $R_{0} = \beta
Tp/(c\delta) > 1$. De factoren: $p/\delta$ is het aantal \omterm{def:b3:virology:virus}{virionen} dat
een besmette cel in haar leven maakt; $\beta T/c$ is het aantal cellen
dat een \omterm{def:b3:virology:virus}{virion} in zijn leven besmet; hun product is het aantal besmette
cellen dat één besmette cel voortbrengt.
\end{solution}

\begin{solution}{exo:b3:virology:11}
$10^{6} = 2^{20}$: twintig halveringstijden, $880$ maanden, ongeveer
$73$ jaar --- langer dan een mensenleven. Middelen die de replicatie
blokkeren raken een provirus dat niet wordt afgeschreven niet aan; het
reservoir blijft bestaan, en zodra de therapie stopt hervat één cel die
weer actief wordt de infectie. Een genezing moet het reservoir
weghalen of blijvend tot zwijgen brengen: de latente cellen doden (ze
onder therapie activeren zodat zij sterven of worden gedood), het
provirus eruit snijden of uitschakelen, of het afweersysteem vervangen
door een dat het \omterm{def:b3:virology:virus}{virus} niet kan besmetten.
\end{solution}

\begin{solution}{exo:b3:virology:12}
$uL = 0.03$ mutaties per genoom; $10^{9}$ genomen per dag dragen
$3\times
10^{7}$ mutaties onder $9\times 10^{4}$ mogelijke enkele mutaties: elk
ontstaat ongeveer $300$ keer per dag. Een bepaalde dubbele mutant
ontstaat met $10^{-12}$ per genoom, $10^{-3}$ per dag --- eens per drie
jaar; een drievoudige met $10^{-18}$, nooit. Twee middelen met
onafhankelijke resistentiemutaties zouden in beginsel volstaan, en drie
geven marge voor slechte therapietrouw. Twee middelen die hetzelfde
enzym binden kunnen door één mutatie worden verslagen die de plaats of
de conformatie van het enzym verandert, en hun resistentiemutaties zijn
niet onafhankelijk; zij tellen als minder dan twee.
\end{solution}

\begin{solution}{pb:b3:virology:1}
\textbf{1.} $2\times 10^{5}\times 1.5\times 10^{4} = 3\times 10^{9}$
\omterm{def:b3:virology:virus}{virionen}.
\textbf{2.} Geklaard $cV = 3\times 3\times 10^{9} \approx 10^{10}$ per
dag; evenveel gemaakt.
\textbf{3.} $I^{*} = cV^{*}/p = 9\times 10^{9}/500 = 1.8\times 10^{7}$
besmette cellen; er sterven er $\delta I^{*} = 9\times 10^{6}$ per dag.
\textbf{4.} RNA $2\times 9700\times 330 = 6.4\times 10^{6}$ Da; eiwit
$2000\times 25\,000 = 5\times 10^{7}$ Da; samen $5.6\times 10^{7}$ Da
$\approx 10^{-16}$ g. Alle \omterm{def:b3:virology:virus}{virionen} samen: $3\times 10^{9}\times
10^{-16} = 3\times 10^{-7}$ g, een derde van een microgram ---
drieduizend keer minder dan een zandkorrel.
\textbf{5.} $1.8\times 10^{7}/10^{12} = 2\times 10^{-5}$ van de
voorraad. Tien jaar van $9\times 10^{6}$ sterfgevallen per dag is
$3\times 10^{10}$ cellen, drie procent van de voorraad: de voorraad
begeeft het niet door een rekenkundige aftrekking maar omdat een
decennium van gedwongen aanvulling en chronische activering het vermogen
om ze te vervangen uitput.
\textbf{6.} $9\times 10^{6}/24 \approx 4\times 10^{5}$ lijken die op
enig moment worden opgeruimd.
\textbf{7.} $V = 2\times 10^{5}(3e^{-0.5t} - 0.5e^{-3t})/2.5$: $t = 1$:
$1.4\times 10^{5}$; $t = 3$: $5.4\times 10^{4}$; $t = 7$: $7\times
10^{3}$; $t = 14$: $2\times 10^{2}$.
\textbf{8.} $2.4\times 10^{5}e^{-0.5t} = 50$: $t = 2\ln(4800) \approx 17$
dagen.
\textbf{9.} De dagen 3--10 liggen op de trage fase: de helling van
$\ln V$ daar is $-\delta$ (van $5.4\times 10^{4}$ naar $1.6\times 10^{3}$
in zeven dagen, $\delta = 0.50$). De snelle fase duurt maar uren
($1/c$), dus vergt $c$ monsters om de paar uur op de eerste dag, aan de
formule met twee exponentiëlen aangepast.
\textbf{10.} De latent besmette cellen, die \omterm{def:b3:virology:virus}{virus} vrijlaten zodra zij
weer actief worden; hun vervalsnelheid is $\ln 2/44$ per maand, ongeveer
$5\times 10^{-4}$ per dag.
\textbf{11.} \omterm{def:b3:virology:virus}{Virion} $\ln 2/3 = \qty{0.23}{d}$, ongeveer \qty{5.5}{h};
besmette cel $\ln 2/0.5 = \qty{1.4}{d}$.
\textbf{12.} Een constante last leek op een \omterm{def:b3:virology:virus}{virus} dat niets deed; de
vergelijkingen tonen haar als een evenwichtstoestand waarin elke dag
$10^{10}$ \omterm{def:b3:virology:virus}{virionen} worden gemaakt en geklaard en $10^{7}$ T-cellen
worden besmet en gedood.
\textbf{13.} $9700\times 10^{-5} \approx 0.1$ mutaties per genoom;
$10^{10}\times 0.1 = 10^{9}$ mutaties per dag.
\textbf{14.} $3\times 9700 \approx \num{29000}$ mogelijke enkele
mutaties; elk ontstaat $10^{9}/\num{29000} \approx 3\times 10^{4}$ keer
per dag.
\textbf{15.} Een bepaalde dubbele: $10^{-10}$ per genoom, ongeveer één
per dag zonder behandeling; een bepaalde drievoudige: $10^{-15}$,
$10^{-5}$ per dag --- eens per drie eeuwen.
\textbf{16.} $10^{7}\times 10^{-15} = 10^{-8}$ per dag, $4\times 10^{-6}$
per jaar.
\textbf{17.} Met één middel weg volstaan mutanten die tegen de andere
twee bestand zijn: een bepaalde dubbele met $10^{-10}$; $7\times 10^{7}$
genomen in de week geven $7\times 10^{-3}$ te verwachten --- nog altijd
onwaarschijnlijk, tenzij de last tijdens de onderbreking naar $10^{10}$
per dag terugveert, waarbij het er één per dag worden en de ontsnapping
waarschijnlijk is. Gemiste doses zijn de manier waarop de resistentie
komt.
\textbf{18.} De analogen delen de actieve plaats van het reverse
transcriptase, en één mutatie daar (of een stel mutaties voor
thymidine-analogen) verandert hoe het enzym met verscheidene ervan
omgaat: kruisresistentie. Twee nucleosiden tellen als minder dan twee
onafhankelijke middelen, zodat de schema's klassen combineren --- een
nucleoside, een integraseremmer, een \omterm{def:b3:virology:drugs}{proteaseremmer} of een
niet-nucleoside.
\textbf{19.} $20$ halveringstijden, $880$ maanden, $73$ jaar; tot
$10^{4}$, $\log_{2}100 = 6.6$ halveringstijden, $290$ maanden, $24$
jaar.
\textbf{20.} Van één \omterm{def:b3:virology:virus}{virion} tot $3\times 10^{9}$ bij $R_{0} = 8$ per
generatie: $\ln(3\times 10^{9})/\ln 8 = 10.5$ generaties, drie weken.
\textbf{21.} Schok en dood: de latente cellen onder therapie weer
activeren zodat zij sterven of worden gedood --- hindernis: geen enkel
middel activeert ze alle, en de geactiveerde cellen worden niet
betrouwbaar gedood. Genbewerking: het provirus eruit snijden of de
receptor CCR5 in de cellen van de patiënt vernietigen --- hindernis:
elke cel bereiken. (Het beenmerg vervangen door donorcellen zonder CCR5
heeft een handvol patiënten genezen, tegen een risico dat alleen
aanvaardbaar is wanneer er toch al een transplantatie nodig was.)
\textbf{22.} $1 - 1/R_{0}$: \qty{67}{\%} voor $R_{0} = 3$, \qty{80}{\%}
voor $5$.
\textbf{23.} De therapie zet $\beta \to 0$ in de patiënt, de last daalt
volgens de stelling tot bijna niets, en de overdracht, die met de last
evenredig is, houdt op; elke besmette persoon behandelen drijft het
$R_{0}$ van de bevolking onder één.
\textbf{24.} $f = (1 - 1/R_{0})/E$: met $R_{0} = 4$ en $E = 0.7$ is $f =
0.75/0.7 = 1.07$ --- meer dan iedereen: het \omterm{def:b3:virology:drugs}{vaccin} alleen kan de drempel
niet halen; met $E = 0.9$ is $f = 0.83$.
\textbf{25.} Ongeveer $10^{10}$ \omterm{def:b3:virology:virus}{virionen} per dag gemaakt; onaantoonbaar
rond dag $17$; zo'n $4\times 10^{-6}$ drievoudige mutanten per jaar
onder therapie.
\end{solution}
